\booksection{Глава 20: Режим разработчика и расширенные права}

\chapterintro{Эта глава предназначена для пользователей, желающих глубже изучить «скрытые функции» телефона. Режим разработчика, Root-права, BL-блокировка... Эти понятия, звучащие как что-то из мира гиков, на самом деле не так далеки от обычного пользователя.}

\booksubsection{I. Режим разработчика: разблокировка скрытых настроек}

Режим разработчика — это встроенное скрытое меню системы телефона, содержащее множество функций, которые не нужны обычным пользователям, но необходимы разработчикам.

\subsubsection*{Как включить:}

\begin{itemize}
  \item \textbf{Android}: Настройки > О телефоне > нажмите 7 раз на «Номер сборки», пока не появится сообщение «Вы стали разработчиком».
  \item \textbf{HarmonyOS}: Настройки > О телефоне > нажмите 7 раз на «Версия HarmonyOS».
  \item \textbf{iOS}: обычный пользователь не может включить полный режим разработчика, требуется платная учётная запись Apple Developer.
\end{itemize}

\subsubsection*{Обзор основных функций:}

\begin{itemize}
  \item \textbf{Отладка по USB}: обязательный переключатель для подключения к компьютеру для расширенных операций (прошивка, восстановление данных, команды ADB). Подробнее см. раздел «Полезные команды ADB» в этой главе.
  \item \textbf{Показывать касания}: отображает точки касания на экране (маленький круг, следующий за пальцем) — удобно для записи обучающих видео или демонстраций.
\end{itemize}

\subsubsection*{Связанное с производительностью и плавностью:}

\begin{itemize}
  \item \textbf{Настройка скорости анимации (Масштаб анимации окна / Масштаб анимации перехода / Длительность анимации)}: эти три параметра управляют длительностью анимации при открытии приложений и переключении страниц. По умолчанию — 1x. Установка на 0,5x делает телефон «визуально» быстрее (фактическая скорость работы программ не меняется, только анимация становится быстрее). Отключение полностью убирает анимацию. Большинство гиков ставят 0,5x — сохраняется плавность переходов, но ощущается быстрее. Полное отключение делает работу резкой и может скрывать реальные задержки программ, создавая ложное ощущение плавности.
  \item \textbf{Принудительная обработка GPU (Force GPU Rendering)}: заставляет всю 2D-графику обрабатываться через GPU, а не CPU. Может улучшить плавность прокрутки на старых моделях или в плохо оптимизированных приложениях, но также увеличивает нагрузку на GPU и расход батареи. Современные телефоны уже широко используют GPU-рендеринг по умолчанию, этот переключатель мало влияет на модели среднего и высокого класса — скорее историческая опция. Обычным пользователям рекомендуется оставить по умолчанию.
  \item \textbf{Профилирование GPU-рендеринга (Profile GPU Rendering)}: отображает на экране в виде гистограммы время рендеринга каждого кадра в реальном времени. Зелёная горизонтальная линия обозначает 16,6 мс (порог времени кадра для 60 fps). Столбики выше зелёной линии означают пропуск кадров. Это мощный инструмент для разработчиков при поиске причин торможения приложений. Обычные пользователи тоже могут использовать его для сравнения, какое приложение оптимизировано хуже — например, если при прокрутке в каком-то приложении столбики часто уходят выше красной линии, то проблема, скорее всего, в нём.
\end{itemize}

\subsubsection*{Связанное с фоном и автономностью:}

\begin{itemize}
  \item \textbf{Ограничение фоновых процессов}: можно установить «Стандартное ограничение», «Не более 1 процесса», «Не более 2 процессов» и т.д. На первый взгляд, ограничение фоновых процессов экономит батарею, но на самом деле встроенный менеджер фона Android уже работает достаточно хорошо. Принудительное ограничение приведёт к частым перезагрузкам приложений (холодный старт, наоборот, потребляет больше энергии и ухудшает опыт). Эта опция больше подходит для отладки. Обычным пользователям не рекомендуется включать её надолго — оставьте «Стандартное ограничение».
  \item \textbf{Не сохранять действия}: каждый раз при уходе с экрана он немедленно уничтожается. Крайне не рекомендуется — приведёт к частым белым экранам и перезагрузкам приложений. Используется только для тестирования логики восстановления приложений.
\end{itemize}

\begin{tcbwarn}{Важные замечания}
Настройки в режиме разработчика могут повлиять на стабильность системы. Не изменяйте непонятные опции. После включения отладки по USB будьте осторожны при подключении к компьютеру — не подключайтесь к ненадёжным компьютерам.

\end{tcbwarn}

\booksubsection{II. OEM-разблокировка: «переключатель прошивки» в параметрах разработчика}

В параметрах разработчика есть малозаметный, но важный переключатель — «OEM-разблокировка» (у некоторых брендов называется «Разрешить OEM-разблокировку» или «Разблокировка устройства»). Сам по себе этот переключатель не разблокирует BL-замок, но это «предварительное разрешение»: только включив его, вы даёте производителю разрешение на разблокировку BL через команды fastboot или фирменные инструменты.

\begin{itemize}
  \item \textbf{Назначение переключателя}: это равносильно заявлению производителю: «Я согласен отказаться от гарантии, принимаю риски на себя, разрешите мне разблокировать BL». Если этот переключатель выключен, любая команда разблокировки BL будет немедленно отклонена.
  \item \textbf{Связь с BL-замком}: OEM-разблокировка — это «подача заявки», а BL-разблокировка — «фактическое действие». Первое делается в настройках телефона одним нажатием, второе требует подачи заявки на сайте производителя и выполнения команд fastboot на компьютере (например,\texttt{fastboot flashing unlock}). Оба шага обязательны.
  \item \textbf{Сброс до заводских настроек сбрасывает?}Да. Переключатель «OEM-разблокировка» после сброса возвращается в выключенное состояние. Но BL-замок после разблокировки сохраняет состояние на физическом уровне — он не запрётся снова после сброса. Для повторной блокировки нужно вручную выполнить\texttt{fastboot flashing lock}.
  \item \textbf{Когда следует включать}: если вы не планируете прошивать телефон или получать Root-права, держите выключенным. Это независимый от «Отладки по USB» переключатель, но при прошивке обычно нужно включать оба.
\end{itemize}

\begin{tcbinfo}{Интересный факт: почему производители добавили этот переключатель}
В ранних версиях Android этого переключателя не было, и вредоносное ПО, получив права, могло тайно прошить систему. Позже Google потребовал от всех производителей Android добавить переключатель «OEM-разблокировка» в параметры разработчика как подтверждение «информированного согласия пользователя». Именно поэтому, даже если вы разблокировали BL-замок, но переключатель OEM-разблокировки выключен, повторная прошивка всё равно не удастся.

\end{tcbinfo}

\booksubsection{III. BL-замок (Bootloader): первые ворота прошивки}

Bootloader — это программа, загружающая операционную систему при запуске телефона. BL-замок — это замок, установленный производителем для безопасности, чтобы предотвратить прошивку неофициальных систем.

\subsubsection*{Функции BL-замка:}

\begin{itemize}
  \item \textbf{Защита безопасности}: предотвращает прошивку вредоносного ПО в систему.
  \item \textbf{Защита от кражи}: даже если телефон перепрошит, без пароля исходной учётной записи он не будет работать.
  \item \textbf{Ограничение гарантии}: разблокировка BL обычно означает потерю официальной гарантии.
\end{itemize}

\subsubsection*{Политика BL-замка у разных брендов:}

\begin{itemize}
  \item \textbf{Xiaomi}: можно разблокировать, требуется привязка аккаунта и период ожидания (7–30 дней).
  \item \textbf{Huawei}: после 2018 года прекратили выдавать коды разблокировки, официальная разблокировка практически невозможна.
  \item \textbf{OnePlus}: можно разблокировать, относительно либеральная политика.
  \item \textbf{OPPO/vivo}: некоторые модели можно разблокировать, требуется подача заявки.
  \item \textbf{Samsung}: можно разблокировать, но это активирует предохранитель Knox, что приведёт к потере Samsung Pay и других функций.
  \item \textbf{Apple}: на устройствах iOS нет традиционного понятия BL-замка, но есть аналогичные механизмы безопасности.
\end{itemize}

\booksubsection{IV. Полезные команды ADB: управление телефоном через компьютер}

ADB — Android Debug Bridge (отладочный мост Android) — это инструмент командной строки, официально предоставляемый Google, позволяющий компьютеру взаимодействовать с Android-телефоном через USB или Wi-Fi. Он не требует Root-прав, достаточно включить «Отладку по USB» на телефоне. Можно представить его как «пульт для подачи команд телефону с компьютера» — прошивка, установка приложений, передача файлов, просмотр логов — всё через него.

\subsubsection*{Как подключиться через USB:}

\begin{enumerate}
  \item На телефоне: войдите в параметры разработчика включите «Отладку по USB».
  \item На компьютере: скачайте Platform Tools с официального сайта Google (содержит исполняемые файлы\texttt{adb}и\texttt{fastboot}), распакуйте в любую папку.
  \item Подключите телефон к компьютеру кабелем. На телефоне появится запрос «Разрешить отладку по USB» — поставьте галочку «Всегда разрешать» и подтвердите.
  \item В командной строке компьютера (PowerShell или cmd в Windows) выполните\texttt{adb devices}. Если вывод похож на\texttt{List of devices attached /\allowbreak  xxxxxxxx device}— подключение установлено.
\end{enumerate}

\begin{tcbwarn}{Не игнорируйте драйверы и запросы авторизации}
Если в Windows команда\texttt{adb devices}показывает «unauthorized» или вообще не показывает устройство, скорее всего, не установлены драйверы или вы забыли нажать «Разрешить» на телефоне. Для Samsung/Xiaomi и других брендов можно установить соответствующий «USB-драйвер» или «Phone Assistant» для установки драйверов. На Mac и Linux драйверы обычно не требуются. Если вы случайно отклонили запрос на телефоне, отключите и снова подключите кабель — запрос появится снова.

\end{tcbwarn}

\subsubsection*{Часто используемые команды ADB:}

\begin{itemize}
  \item \texttt{adb install app.\allowbreak apk}: установить APK с компьютера на телефон (\texttt{-\allowbreak r}— переустановка с сохранением данных;\texttt{-\allowbreak s}— установка на SD-карту).
  \item \texttt{adb uninstall com.\allowbreak xxx.\allowbreak xxx}: удалить приложение по имени пакета.
  \item \texttt{adb push путь\_на\_ПК /\allowbreak sdcard/\allowbreak путь\_назначения}: отправить файл с компьютера на телефон.
  \item \texttt{adb pull /\allowbreak sdcard/\allowbreak файл путь\_на\_ПК}: скачать файл с телефона на компьютер (гораздо быстрее, чем отправлять себе в WeChat и потом скачивать).
  \item \texttt{adb shell}: войти в Linux Shell телефона, можно выполнять команды\texttt{ls},\texttt{pm},\texttt{am}и т.д.
  \item \texttt{adb reboot bootloader}: перезагрузить телефон в режим Fastboot (вход для прошивки).
  \item \texttt{adb reboot recovery}: перезагрузить в режим Recovery.
  \item \texttt{adb shell pm list packages}: вывести список всех установленных приложений с именами пакетов (обязательно для поиска имени пакета). Добавьте\texttt{-\allowbreak s}для просмотра только системных,\texttt{-\allowbreak 3}— только сторонних.
  \item \texttt{adb shell pm uninstall -\allowbreak -\allowbreak user 0 имя\_пакета}: удалить предустановленное приложение для текущего пользователя без Root-прав — подробнее ниже.
\end{itemize}

\subsubsection*{Практика: удаление предустановленных приложений без Root:}

Это одна из самых полезных функций ADB: команда\texttt{pm uninstall -\allowbreak -\allowbreak user 0}позволяет удалить предустановленные приложения для «текущего пользователя» без Root-прав и без изменения системного раздела. Удалённое приложение «исчезает» из системы, но физически остаётся в системном разделе — после сброса до заводских настроек оно вернётся.

\begin{enumerate}
  \item Сначала командой\texttt{adb shell pm list packages -\allowbreak s}выведите список всех системных приложений и найдите имя пакета для удаления. Например, браузер Xiaomi —\texttt{com.\allowbreak android.\allowbreak browser}, магазин приложений Huawei —\texttt{com.\allowbreak vmall.\allowbreak client}.
  \item Не уверены, какому приложению соответствует имя пакета? На компьютере используйте\texttt{adb shell pm list packages | findstr ключевое\_слово}для поиска. Или установите стороннее приложение «Package Viewer» для просмотра имён пакетов прямо на телефоне.
  \item Выполните удаление:\texttt{adb shell pm uninstall -\allowbreak -\allowbreak user 0 имя\_пакета}, например\texttt{adb shell pm uninstall -\allowbreak -\allowbreak user 0 com.\allowbreak android.\allowbreak browser}. Ответ\texttt{Success}означает успех.
  \item Как восстановить? Два способа: (1) Сброс до заводских настроек (самый надёжный, но трудоёмкий); (2) Переустановка:\texttt{adb shell cmd package install-\allowbreak existing имя\_пакета}.
\end{enumerate}

\textbf{Ни в коем случае не удаляйте ключевые системные пакеты}\par

Хотя «pm uninstall --user 0» — безопасный способ без Root, удаление критических системных приложений (таких как\texttt{com.\allowbreak android.\allowbreak systemui}— строка состояния,\texttt{com.\allowbreak android.\allowbreak phone}— телефон,\texttt{com.\allowbreak android.\allowbreak settings}— настройки) приведёт к краху системы, чёрному экрану, невозможности загрузки и другим серьёзным проблемам. Перед удалением погуглите «можно ли удалить этот пакет». Лучше удалить меньше, чем удалить лишнее. Рекомендуется удалять по одному-два за раз, убедиться, что проблем нет, и продолжать.

\begin{tcbinfo}{Интересный факт: что означает --user 0}
Начиная с Android 4.2 поддерживается многопользовательский режим, у каждого пользователя свой ID (владелец — 0, дополнительные пользователи начинаются с 10).\texttt{-\allowbreak -\allowbreak user 0}означает «удалить только для текущего владельца». Файл APK в системном разделе остаётся, просто не показывается владельцу. Именно поэтому это не требует Root — системные файлы не затрагиваются, изменяется только «конфигурация пользователя». В этом и заключается главное отличие от удаления с Root-правами: Root-удаление действительно удаляет из системы для всех пользователей; ADB-удаление — только «скрывает для текущего пользователя».

\end{tcbinfo}

\booksubsection{V. Root-права: «суперадминистратор» Android}

Root-права эквивалентны правам суперпользователя в Linux (полный контроль над корневым каталогом /, возможность чтения и записи любых файлов) — это высший уровень прав в системе Android. Они принципиально отличаются от Администратора в Windows: Windows из соображений безопасности накладывает множество ограничений на администратора (настоящие высшие права принадлежат TrustedInstaller), и даже Administrator не может, как root в Linux, напрямую выполнить\texttt{sudo rm -\allowbreak rf /\allowbreak *}. Получив Root-права, вы можете изменять системные файлы, удалять предустановленные приложения, использовать продвинутые приложения, требующие Root.

\subsubsection*{Преимущества Root:}

\begin{itemize}
  \item Удаление неудаляемых предустановленных приложений.
  \item Использование продвинутых приложений (например, Titanium Backup, Greenify).
  \item Изменение системного интерфейса и функций.
  \item Установка сторонних Recovery и кастомных ROM.
\end{itemize}

\subsubsection*{Риски Root:}

\begin{itemize}
  \item \textbf{Потеря гарантии}: большинство производителей не обслуживают телефоны с Root.
  \item \textbf{Риски безопасности}: вредоносные приложения могут получить Root-права с серьёзными последствиями.
  \item \textbf{Нестабильность системы}: ошибочные действия могут привести к краху системы.
  \item \textbf{Банковские приложения не работают}: многие банковские и платёжные приложения отказываются работать при обнаружении Root.
  \item \textbf{Сбой OTA-обновлений}: возможна невозможность получения официальных обновлений системы.
\end{itemize}

\textbf{Серьёзное предупреждение}\par

Операции с Root сопряжены с определённым риском и могут превратить телефон в «кирпич» (невозможность включения). Если вы новичок, рекомендуется сначала тщательно изучить соответствующую информацию и убедиться в наличии полной резервной копии. Перед операцией внимательно прочитайте инструкцию иподтвердите вашу модель и версию системы.

\booksubsection{VI. SafetyNet и Play Integrity: почему после Root не работают банковские приложения}

Многие пользователи с Root сталкивались с этим моментом краха: после Root с энтузиазмом установили кучу модулей, а потом открыли банковское приложение — «Ваше устройство небезопасно, использование невозможно». Google Pay тоже не работает, даже Netflix отказывается от HD-воспроизведения. За этим стоят два механизма проверки целостности устройства от Google: старый SafetyNet и новый Play Integrity API.

\subsubsection*{Что такое SafetyNet / Play Integrity:}

\begin{itemize}
  \item \textbf{SafetyNet Attestation (старая версия, с 2017)}: ранний API проверки целостности устройства от Google. Приложения через него запрашивают у серверов Google, «проходит ли устройство тест совместимости CTS и находится ли в неизменённом официальном состоянии». Результаты делятся на Basic Attestation (базовая проверка) и Strong Attestation (строгая проверка, зависит от аппаратного ключа). Google объявил о постепенном отказе от SafetyNet после 2024 года.
  \item \textbf{Play Integrity API (новая версия, с 2022)}: замена от Google, более строгая и всеобъемлющая, чем SafetyNet. Помимо целостности устройства, включает «целостность приложения» (не было ли оно изменено) и «целостность аккаунта/авторизации» (доверяет ли Google аккаунту пользователя). Новые приложения должны переходить на Play Integrity.
  \item \textbf{Кто использует}: банки, финансы, платёжные приложения (Google Pay, мобильные банки, часть сценариев Alipay/WeChat Pay), стриминг (Netflix — защита авторских прав HD), некоторые игры с античитом (например, игры Tencent/NetEase), корпоративные приложения (MDM — проверка соответствия устройства).
\end{itemize}

\subsubsection*{Три уровня сертификации Play Integrity:}

\begin{itemize}
  \item \textbf{MEETS\_DEVICE\_INTEGRITY (целостность устройства пройдена)}: устройство работает на подлинной системе Android, Bootloader не разблокирован, Root отсутствует, тест совместимости CTS пройден. Можно нормально использовать банковские и платёжные приложения.
  \item \textbf{MEETS\_BASIC\_INTEGRITY (базовая целостность пройдена)}: устройство в основном надёжно, но могут быть незначительные аномалии. Некоторые приложения будут работать в ограниченном режиме (Netflix может воспроизводить, но с ограничением разрешения). Некоторые строгие банковские приложения откажут напрямую.
  \item \textbf{MEETS\_STRONG\_INTEGRITY (строгая целостность пройдена)}: высший уровень, требует подтверждения аппаратным ключом (например, чип Titan M в серии Google Pixel). Большинство моделей сторонних брендов даже без Root не достигают этого уровня.
  \item \textbf{Без каких-либо меток (NO\_INTEGRITY)}: устройство имеет Root, разблокированный BL, прошито сторонней ROM или распознано как эмулятор. Подавляющее большинство платёжных и банковских приложений откажут в работе.
\end{itemize}

\subsubsection*{Почему эти приложения перестают работать после Root:}

Приложения через Play Integrity API отправляют запрос на серверы Google, которые проверяют множество показателей устройства — состояние блокировки Bootloader, изменён ли системный раздел, существует ли бинарный файл su, находится ли SELinux в разрешающем режиме и т.д. Любая аномалия приводит к снижению уровня сертификации или его отсутствию. Банковское приложение видит «устройство небезопасно» и, естественно, отказывает в обслуживании.

\subsubsection*{Современные методы скрытия: Magisk Hide / Shamiko / Zygisk:}

Чтобы пользоваться Root-правами и одновременно банковскими приложениями, сообщество разработало ряд решений для «скрытия состояния Root». Основная идея: сделать так, чтобы проверяющая программа видела не изменённую систему, а «выглядящее нормально» устройство.

\begin{itemize}
  \item \textbf{Magisk}: самый популярный на сегодня метод Systemless Root — не изменяет системный раздел, а реализует Root через монтирование. Такой подход снижает вероятность сбоя OTA-обновлений и позволяет относительно чисто удалить Root. Это основа для всех последующих методов скрытия. Однако для использования Magisk всё равно требуется разблокировка BL-замка.
  \item \textbf{Magisk Hide (разработка прекращена)}: ранний встроенный модуль Magisk для скрытия Root от указанных приложений. С 2021 года официально удалён из Magisk, так как Google постоянно обновляет методы обнаружения — это превратилось в «игру в кошки-мышки».
  \item \textbf{Zygisk}: новый механизм, введённый в Magisk 24.0+, аналогичный Riru. Загружает модули через инъекцию в процесс Zygote, что скрытнее прямой модификации системы. Включение Zygisk — предварительное условие для использования современных методов скрытия.
  \item \textbf{Shamiko}: самый популярный на сегодня модуль скрытия, работает в паре с Zygisk. Может скрывать Root, сам Magisk и подделывать состояние блокировки Bootloader для указанных приложений (режим «обратного белого списка»). Shamiko распространяется через репозиторий LSPosed или GitHub.
  \item \textbf{Play Integrity Fix (osm0sis)}: модуль Magisk, специально нацеленный на обнаружение Google Play Integrity. Изменяет отпечаток устройства (fingerprint), чтобы Google считал ваше устройство одной из официальных моделей без Root.
\end{itemize}

\begin{tcbwarn}{Игра в кошки-мышки никогда не закончится}
Google периодически обновляет правила обнаружения, а разработчики приложений следуют за обновлениями проверок Play Integrity. То, что работает сегодня, завтра может быть распознано. Если вы полностью зависите от банковских и платёжных приложений (ежедневные переводы, оплата по QR-коду), самый надёжный выбор — не делать Root и не разблокировать BL-замок. Методы скрытия больше подходят для конфигурации «основной телефон без Root + запасной телефон с Root для экспериментов».

\end{tcbwarn}

\booksubsection{VII. Рекомендации для обычных пользователей}

\begin{itemize}
  \item \textbf{Режим разработчика}: можно включить, попробовать ускорение анимации, профилирование GPU и другие функции, но не изменяйте непонятные опции.
  \item \textbf{OEM-разблокировка и BL-замок}: если вы не планируете прошивку или Root, безопаснее всего держать оба переключателя в состоянии по умолчанию (выключено).
  \item \textbf{ADB}: мощный инструмент без Root, доступный обычным пользователям — например, удаление предустановленных приложений командой\texttt{pm uninstall -\allowbreak -\allowbreak user 0}. Но перед удалением обязательно проверьте, можно ли удалять этот пакет.
  \item \textbf{Root}: обычным пользователям Root не рекомендуется — риски превышают выгоду. Если вам действительно нужны определённые функции, сначала узнайте, есть ли альтернативы без Root (например, удаление предустановленных приложений через ADB, использование Shizuku для замены части функций Root).
\end{itemize}

\begin{tcbpractice}{Попробуйте: включите режим разработчика и настройте скорость анимации}
\begin{enumerate}
  \item Зайдите в Настройки > О телефоне.
  \item Нажмите 7 раз на «Номер сборки», пока не появится сообщение об активации режима разработчика.
  \item Вернитесь в Настройки, найдите «Параметры разработчика» (обычно внизу настроек или в разделе «Система»).
  \item Найдите «Масштаб анимации окна», «Масштаб анимации перехода», «Длительность анимации» и установите все на 0,5x или отключите.
  \item Почувствуйте, стал ли телефон быстрее.
\end{enumerate}

\end{tcbpractice}
